\section{Proposed Method}
\label{sec:method}

\subsection{Experimental Design}

This study employs a controlled within-subject design in which every benchmark task is attempted under both experimental conditions. Condition~A (the baseline) represents the conventional flat tool-use paradigm, while Condition~B (the experimental intervention) implements progressive skill disclosure via the OpenClaw \cite{openclaw_2025} framework. All other variables---the underlying language models, the task inputs, the skill library content, and the evaluation criteria---are held constant across both conditions. This design ensures that any observed differences in outcome accuracy, trajectory quality, or token efficiency are attributable exclusively to the capability-loading mechanism rather than to differences in model capability or task difficulty.

\subsection{Technology Stack}

All experiments are implemented in Python using the following tools. No additional agent orchestration frameworks beyond OpenClaw are required.

\begin{itemize}
    \item \textbf{Python 3.11:} Task dispatch, action logging, metric computation, and result storage. All measurement values are written to structured JSON files for reproducibility.
    \item \textbf{Ollama} \cite{ollama_2024}: Local inference server for small language models. Ollama handles model download, quantisation, and GPU or CPU scheduling automatically, providing a consistent HTTP API across all local model backends.
    \item \textbf{Gemini 3.1 Pro} \cite{google_gemini_2025} \textit{(primary cloud baseline):} Accessed via the Google Generative AI Python SDK. Token counts and latency are retrieved from the API response metadata.
    \item \textbf{GPT-4o} \cite{openai_gpt4o_2024} \textit{(secondary cloud baseline):} Accessed via the OpenAI Python SDK. Token billing records are used for cost calculation.
    \item \textbf{OpenClaw} \cite{openclaw_2025} \textit{(Condition B only):} Provides the native three-tier skill architecture. Python intercepts the agent's \texttt{load\_skill} actions to log each disclosure event and measure the token delta introduced by loading each SKILL.md file.
\end{itemize}

\subsection{Agent Configurations}

\subsubsection*{Condition A --- Flat Tool Use (Baseline)}

In Condition~A, the complete library of 20 available skills is converted to a flat JSON schema list and injected into the agent's system prompt at the beginning of every task. Each entry contains the skill name, a one-sentence description, and the input and output parameter types. The agent reasons over this list at every step and selects the appropriate function by name. This configuration mirrors the standard function-calling implementation offered by the OpenAI API \cite{openai_gpt4o_2024} and replicated identically in the Python controller for local model backends.

\subsubsection*{Condition B --- Progressive Skill Disclosure (Experimental)}

In Condition~B, the agent is initialised with only a compact metadata index: a numbered list of skill names paired with one-line descriptions, occupying fewer than 500 tokens. When the agent determines that a skill is required, it emits a \texttt{load\_skill(name)} action. The Python controller intercepts this action, reads the corresponding SKILL.md file from disk, and appends its full content---including applicability conditions, execution policy, and termination criteria---to the active message thread. Once the skill's termination criteria are satisfied, the instructional content is removed from the thread before the subsequent reasoning step, restoring the context to the compact metadata state. The executable scripts associated with a skill are invoked directly through Python subprocess calls and are never loaded into the language model's context.

\subsection{Language Model Backends}

Three language model backends are evaluated under both conditions to assess whether the performance gap between paradigms generalises across model capability levels:

\begin{enumerate}
    \item \textbf{\texttt{gemma4:4b}} (Google DeepMind, 2026) \cite{gemma4_2026}: A 4.5-billion-parameter multimodal model served locally via Ollama. Selected as the primary local baseline owing to its strong document understanding capability and efficient on-device inference.
    \item \textbf{\texttt{qwen3.5:4b}} (Alibaba, 2025) \cite{qwen3_2025}: A 4-billion-parameter model with demonstrated strength in structured information extraction, also served via Ollama.
    \item \textbf{Gemini 3.1 Pro} (Google DeepMind) \cite{google_gemini_2025}: The primary cloud-based frontier model, accessed via the Google Generative AI SDK. Serves as the upper-bound performance reference and is expected to exhibit less sensitivity to context saturation given its substantially larger context window.
    \item \textbf{GPT-4o} (OpenAI) \cite{openai_gpt4o_2024}: The secondary cloud-based baseline, included to assess whether findings are consistent across cloud providers and model families.
\end{enumerate}

\subsection{Skill Library}

A controlled library of 20 skills is constructed for this study, organised into three functional domains:

\begin{enumerate}
    \item \textbf{Document Processing (8 skills):} Extracting named fields from unstructured PDFs and images, summarising document content, converting between document formats, and validating extracted data against a schema.
    \item \textbf{Web and Data Operations (6 skills):} Performing structured web searches, fetching and parsing remote content, normalising tabular data, and formatting output as structured JSON or CSV.
    \item \textbf{File Management (6 skills):} Reading from and writing to local files, organising directories, comparing file versions, and generating formatted reports.
\end{enumerate}

Each SKILL.md file contains the three mandatory structural components defined by the progressive disclosure architecture: (1) \textit{Applicability Conditions}, specifying the environmental preconditions required before the skill may be invoked; (2) \textit{Execution Policy}, a numbered sequence of procedural steps with explicit error-handling branches; and (3) \textit{Termination Criteria}, defining the conditions that indicate successful completion or grounds for aborting the operation. In Condition~A, the same 20 skills are represented as JSON schema objects loaded into the system prompt. Two library sizes are evaluated---10 skills and 20 skills---to measure how the performance gap between conditions scales with the number of available capabilities.

\subsection{Task Benchmark}

A benchmark of 40 tasks is constructed across three difficulty levels:

\begin{itemize}
    \item \textbf{Level 1 --- Single-skill tasks (15 tasks):} Each task requires the identification and execution of exactly one skill. Example: \textit{``Extract the invoice number, vendor name, and total amount due from the attached PDF.''}
    \item \textbf{Level 2 --- Two-skill sequential tasks (15 tasks):} Each task requires two skills applied in a defined sequence. Example: \textit{``Extract the supplier contact details from this form and write them to a formatted CSV file.''}
    \item \textbf{Level 3 --- Multi-skill compositional tasks (10 tasks):} Each task requires three or more skills with at least one conditional branch based on intermediate results. Example: \textit{``Search for the most recent pricing data for the three vendors named in this document, validate the figures against the attached budget table, and produce a discrepancy report.''}
\end{itemize}

All 40 tasks are executed under both conditions and across all four model backends. Each task is repeated three times per condition using a fixed random seed to allow computation of mean and standard deviation values, yielding a total of $40 \times 2 \times 4 \times 3 = 960$ individual agent runs.

\subsection{Measures}

\subsubsection*{Outcome Metrics}

\begin{itemize}
    \item \textbf{Task Completion Rate:} The proportion of tasks for which the agent produces a correct final output, scored using the ANLS* metric \cite{anlsstar_2024} for text-extraction tasks and binary pass/fail scoring for file and web tasks.
\end{itemize}

\subsubsection*{Trajectory Metrics}

\begin{itemize}
    \item \textbf{Hallucinated Skill Invocations:} The count of instances in which the agent attempts to invoke a skill name that does not exist in the library, or invokes an existing skill with semantically incorrect parameters that violate its applicability conditions.
    \item \textbf{Trajectory Compliance Score:} A per-task score measuring the degree to which the agent followed the numbered steps of the relevant execution policy in the correct order, computed by comparing the logged action trace against the ground-truth policy graph for each skill.
    \item \textbf{Prompt Token Count:} The total number of prompt tokens consumed across all reasoning steps for a given task, obtained from Ollama server logs for local models and from API response metadata for cloud models.
    \item \textbf{Error Recovery Rate:} The proportion of tasks in which the agent encountered a tool execution failure and autonomously recovered within three subsequent reasoning steps.
\end{itemize}

\subsubsection*{Efficiency Metrics}

\begin{itemize}
    \item \textbf{End-to-End Latency:} Wall-clock time from task submission to final output, measured in seconds and averaged across all repetitions within each task level.
    \item \textbf{Inference Cost:} For cloud models, the monetary cost derived from token billing records (number of tokens multiplied by the provider's published per-million-token rate). For local models, the energy cost is calculated by multiplying the measured average wattage---recorded using \texttt{powermetrics} on macOS or \texttt{powertop} on Linux \cite{tokenpowerbench_2024}---by elapsed time and the local electricity rate, supplemented by a hardware depreciation share calculated using straight-line depreciation over three years.
\end{itemize}
